\documentclass[12pt]{article}
\usepackage{graphicx} % Required for inserting images
\usepackage[margin=0.5in]{geometry}
\usepackage{amsmath, amssymb}
\usepackage{mathrsfs}


\usepackage{soul}


% Dr. David Brown Commands
\newcommand{\ds}{\displaystyle}
\newcommand{\R}{\mathbb{R}}
\newcommand{\N}{\mathbb{N}}
\newcommand{\Q}{\mathbb{Q}}
\newcommand{\C}{\mathbb{C}}


% Commands to get script r
\usepackage{calligra}
\DeclareMathAlphabet{\mathcalligra}{T1}{calligra}{m}{n}
\DeclareFontShape{T1}{calligra}{m}{n}{<->s*[2.2]callig15}{}
\newcommand{\scripty}[1]{\ensuremath{\mathcalligra{#1}}}

% Unit commands
\newcommand{\degree}{\text{°}}
\newcommand{\unitSpeed}{\frac{\text{m}}{\text{s}}}
\newcommand{\unitAcceleration}{\frac{\text{m}}{\text{s}^2}}

% Constant commands
\newcommand{\avogadro}{6.022\times10^{23}}

\newcommand{\boltzmannConstK}{1.381\times10^{-23}\frac{\text{J}}{\text{K}}}
\newcommand{\boltzmannConstInEV}{8.617\times10^{-5}\frac{\text{eV}}{\text{K}}}
\newcommand{\planckConst}{6.626\times10^{-34}\text{J}\cdot\text{s}}

\newcommand{\atmP}{1.01325\times10^5\,\text{Pa}}

% Commands to easily write partial derivatives
\newcommand{\partDeriv}[2]{\frac{\partial #1}{\partial #2}}
\newcommand{\doublePartDeriv}[2]{\frac{\partial^2 #1}{\partial^2 #2}}


\title{Problem 6.41}
\author{Gabriel Decker}


\begin{document}



\maketitle
\begin{flushleft}


For this problem, we will follow a similar process as was described in the book for the three-dimensional case of the velocity distribution function $\mathcal{D}$.
We will logically assume that $\mathcal{D}$ is proportional to the probability of a molecule having a particular velocity times the number of velocities that correspond to a speed.\\~\\

Recall that according to the Boltzmann distribution, probability is related to energy by the following.
\begin{equation}
    \mathcal{P}(s)\propto e^{-E(s)/kT}
\end{equation}
In this case, we can use the kinetic energy equation.
\begin{equation}
    E=\frac{1}{2}mv^2
\end{equation}
where $v=|\vec{v}|$.
Therefore, we can combine these into the following.
\begin{equation}
    \mathcal{P}\propto e^{-mv^2/2kT}
\end{equation}\\~\\

When we integrate our eventual distribution function, it will be integrated over speed, not velocity.
The above proportionality considers the velocity, not the speed, so we must convert it from velocity space to speed space.\\~\\

The result of the above work is the following distribution function based on velocity.
\begin{equation}
    \mathcal{D}_{\vec{v}}=C\cdot e^{-mv^2/2kT}
\end{equation}
We want a distribution function that gives the probability of a particle having a velocity between $v$ and $v+dv$.
This would be given by $\mathcal{D}(v)\,dv$.
We want to convert $\mathcal{D}_{\vec{v}}$ into $\mathcal{D}(v)\,dv$.
To convert it to distribution that only depends on the speed, we must integrate it over velocity space.
This will be done in polar corrdinates with $da=v\,dv\,d\theta$.
$$\mathcal{D}(v)\,dv=\left[\int_0^{2\pi}\mathcal{D}(v)_{\vec{v}}\cdot v\,d\theta\right]\,dv$$
$$\implies\mathcal{D}(v)\,dv=2\pi v\mathcal{D}(v)_{\vec{v}}\,dv$$
$$\implies\mathcal{D}(v)=2\pi v\mathcal{D}(v)_{\vec{v}}$$
$$\implies\mathcal{D}(v)=2\pi vCe^{-mv^2/2kT}$$\\~\\

Solving for $\mathcal{D}$ implies finding $C$. This can be done by assuming that the probability of finding the particle at some speed is 100\%.
$$\mathcal{P}(v=[0,\infty])=1=2\pi\int_0^\infty v\mathcal{D}(v)_{\vec{v}}\,dv$$
$$\implies1=2\pi\int_0^\infty Cve^{-mv^2/2kT}\,dv$$\\~\\

To solve this, we'll do a change of variables with $x=v\sqrt{m/2kT}$ or $dx=\sqrt{m/2kT}\,dv$.
$$\implies1=2\pi C\cdot\frac{2kT}{m}\int_0^\infty xe^{-x^2}\,dx$$
$$\implies1=C\cdot\frac{4\pi kT}{m}\int_0^\infty xe^{-x^2}\,dx$$
By using the sympy code found in the included script, we get the following.
$$\implies1=C\cdot\frac{2\pi kT}{m}$$
$$\implies C=\frac{m}{2\pi kT}$$\\~\\

Plugging in our $C$ into our $\mathcal{D}$ equation, we can finalize our velocity distribution function.
$$\mathcal{D}(v)=2\pi vCe^{-mv^2/2kT}$$
$$\implies\mathcal{D}(v)=\frac{mv}{kT}e^{-mv^2/2kT}$$









\end{flushleft}

\end{document}
